
% Place this part in the finite-record section near the operational estimator.
\subsection{True-law distance and finite-record confidence sets}
\label{subsec:empirical-topology}

The observation topology in the stability theorems is a distance between
\emph{true} row laws. It is not the total-variation estimation error of a raw
empirical measure. For every finite nonempty row,
$\widehat\mu_I=n_I^{-1}\sum_{m:I_m=I}\delta_{(\tau_m,J_m)}$ is atomic, while
the CTMC row law has a density in time. The finite set of observed atoms has
empirical probability one and true probability zero, so
\begin{equation}
 \TV(\widehat\mu_I,\mu_I)=1.\label{eq:empirical-tv-one}
\end{equation}
Weak convergence and convergence of bounded Laplace tests remain possible.
The moment estimator above uses precisely those bounded tests and therefore
does not assume a shrinking empirical-TV error. A smoothing or detector model
would need its own error analysis.

An elementary confidence rectangle gives a finite-sample alternative to
the pointwise delta approximation.
\begin{proposition}[Sequential moment confidence rectangle]
\label{prop:moment-rectangle}
Fix $K$ positive Laplace arguments in advance, $0<\alpha<1$, and a deterministic
number $n\ge1$ of complete inter-mark intervals. Let $d_{\rm mom}=4K$ and
\begin{equation}
 t_{n,\alpha}=\sqrt{\frac n2\log\frac{2d_{\rm mom}}\alpha}.
\end{equation}
For the empirical joint moments~\eqref{eq:empirical-laplace},
\begin{equation}
 \mathbb P_Q\left\{
 \left|\widehat\mu_{I,jL}-\widehat\Psi_{Q,IL}(\lambda_j)\right|
 \le\frac{t_{n,\alpha}}{n_I}
 \ \text{for all }(I,j,L)\text{ with }n_I>0\right\}\ge1-\alpha.
 \label{eq:finite-moment-rectangle}
\end{equation}
Rows with $n_I=0$ impose no constraint. Stationarity of the initial mark is
unnecessary; the exact CTMC and resolved-reset observation assumptions suffice.
\end{proposition}
\begin{proof}
Let $\mathcal F_0$ contain the first observed mark and let $\mathcal F_m$
include the next $m$ complete marked intervals. Then $I_m$ is
$\mathcal F_{m-1}$-measurable. For one fixed coordinate define
\begin{equation}
 Z_m=\mathbf1\{I_m=I\}
 \left(e^{-\lambda_j\tau_m}\mathbf1\{J_m=L\}
       -\widehat\Psi_{Q,IL}(\lambda_j)\right).
\end{equation}
The reset property gives $\mathbb E[Z_m\mid\mathcal F_{m-1}]=0$.
Conditionally, $Z_m$ lies in an interval of length at most one, so
conditional Hoeffding's lemma yields
$\mathbb E[e^{\eta Z_m}\mid\mathcal F_{m-1}]\le e^{\eta^2/8}$.
Iteration and Chernoff's bound, optimized at $\eta=4t/n$, imply
\begin{equation}
 \mathbb P_Q\left\{\left|\sum_{m=1}^nZ_m\right|>t\right\}
 \le2e^{-2t^2/n}.
\end{equation}
Union over $d_{\rm mom}$ coordinates and use
$\sum_mZ_m=n_I(\widehat\mu_{I,jL}-\widehat\Psi_{Q,IL}(\lambda_j))$
for nonempty rows. The argument neither assumes independent successive rows
nor conditions on their random counts.
\end{proof}

Intersecting the simultaneous rectangle with a declared microscopic model
class gives a confidence set for the generator with coverage at least
$1-\alpha$. Its full entropy image, or a certified enclosing interval,
inherits this coverage and may have an infinite upper endpoint. Finding a
large feasible entropy by numerical optimization does not certify that
endpoint from above. The bound is conservative and does not establish
efficient inversion or an efficient estimator. It applies to fixed $n$
complete intervals and preselected arguments; an adaptive stopping rule,
same-record design selection, terminal censoring at fixed physical time,
or missed events requires additional analysis.

The existing simulations already use actual sequential CTMC records and
dependence-aware moment covariance. Their reported bias, variance, and
Monte Carlo coverage remain operationally meaningful. Their Wald intervals
and low-count warning retain the stated pointwise or heuristic status;
Eq.~\eqref{eq:finite-moment-rectangle} is a separate finite-sample
confidence-set construction, not a retroactive coverage guarantee for those
intervals.
